



Orbiter can create objects in projective space.
To do so, define an object of type
\verb'-geometric_object'.
The definition of a geometric object requires a projective geometry object. 
For this reason, the definition requires an extra argument, 
which is the label of a previously created projective geometry object.
After that, one of the commands shown in 
Tables~\ref{tab:geometric:objects:1} and~\ref{tab:geometric:objects:2} can be issued.
\orbitertableofcommands{geometric_object_1}
\orbitertableofcommands{geometric_object_2}




The following command creates an elliptic quadric ovoid on $\PG(3,8)$:

\sourcecode{elliptic_quadric_ovoid_q8}

The next command creates the Suzuki-Tits ovoid in $\PG(3,8)$:

\sourcecode{ovoid_ST_q8}




\bigskip

The following command creates the Baer subplane of $\PG(2,4)$:

\sourcecode{Baer_PG_2_4}

The next command creates the Baer subgeometry of $\PG(3,4)$:

\sourcecode{Baer_PG_3_4}


\bigskip

Let us consider BLT-sets. The following command recalls the unique BLT-set 
of order $5$ from the database:

\sourcecode{BLT_database_5_0}

Let us pull the first BLT-set of order $7$ and print it to latex:

\sourcecode{BLT_database_7_0_print}

The output is shown below:
\orbiteroutput{GRAPHICS/LATEX/BLT_7_0}


\bigskip


